\documentclass[12pt, a4paper]{article}
\usepackage{geometry}
\geometry{margin=1in}
\usepackage{amsmath, amssymb}
\usepackage{hyperref}
\usepackage{listings}
\usepackage{xcolor}
\usepackage{titlesec}
\usepackage{enumitem}
\usepackage{booktabs}
\usepackage{array}

\hypersetup{
    colorlinks=true,
    linkcolor=blue,
    urlcolor=cyan,
}

\titleformat{\section}{\large\bfseries}{}{0em}{}[\titlerule]
\titleformat{\subsection}{\normalsize\bfseries}{}{0em}{}

\title{\textbf{Blockchain: A Ground-Up Mastery Guide}\\
\large Built via Feynman Method for CeDAR (LUMS) Preparation}
\author{Shaheer}
\date{\today}

\begin{document}

\maketitle
\tableofcontents
\newpage

% ============================================================
\section{Phase 1: Foundations}
% ============================================================

\subsection{The Core Problem Blockchain Solves}

Every financial system in history has relied on a \textbf{trusted middleman} --- a bank,
a government, a court --- to maintain a record of who owns what. When you send money,
you are not moving anything physical; you are asking your bank to update a number in
\emph{their private database}.

The fundamental vulnerabilities of this model:
\begin{itemize}
    \item The middleman can be corrupt or coerced.
    \item The middleman is a single point of failure.
    \item Access can be denied or frozen.
    \item Cross-border trust is fragile.
\end{itemize}

\textbf{Satoshi Nakamoto's insight (2008):} Replace the single trusted middleman with
\emph{thousands of strangers} each holding an identical copy of the same ledger, all
of whom must agree before any change is accepted. This is the blockchain.

\begin{center}
\begin{tabular}{>{\bfseries}l l}
\toprule
Traditional System & Blockchain \\
\midrule
One private database & Thousands of identical copies \\
Trust the institution & Trust the math \\
Single point of failure & No central point to attack \\
Permission required & Open and permissionless \\
\bottomrule
\end{tabular}
\end{center}

% ------------------------------------------------------------
\subsection{Why ``Chain''? The Structure of a Blockchain}

A blockchain is composed of \textbf{blocks} linked in a chain. Each block contains:
\begin{enumerate}
    \item A set of \textbf{transactions} (the data).
    \item The \textbf{hash of the previous block} --- this is the ``chain'' link.
    \item Its own \textbf{hash} (a fingerprint of its entire contents).
\end{enumerate}

\begin{verbatim}
┌──────────────────────┐     ┌──────────────────────┐
│  BLOCK #2            │     │  BLOCK #3            │
│  Txns: A→B: 500      │     │  Txns: B→C: 200      │
│  Prev Hash: [Block1] │◄────│  Prev Hash: [Block2] │
│  Own Hash:  a3f2c7.. │     │  Own Hash:  7b84fa.. │
└──────────────────────┘     └──────────────────────┘
\end{verbatim}

\textbf{Why chaining prevents tampering:} If an attacker modifies any transaction
in Block 2, Block 2's hash changes entirely. This invalidates Block 3's
``Prev Hash'' reference, which cascades and breaks every subsequent block.
The fraud becomes immediately visible to all participants.

Two distinct ideas work together:

\begin{center}
\begin{tabular}{>{\bfseries}l l}
\toprule
Mechanism & What It Prevents \\
\midrule
Chaining (linked hashes) & Secretly editing old history \\
Distribution (thousands of copies) & Any single party controlling the record \\
\bottomrule
\end{tabular}
\end{center}

% ------------------------------------------------------------
\subsection{The Hash Function: The Mathematical Fingerprint}

A \textbf{hash function} takes any input and produces a fixed-size output (the hash).
Bitcoin uses \textbf{SHA-256}, which always produces a 256-bit output.

\textbf{The four critical properties:}

\begin{enumerate}
    \item \textbf{Deterministic:} The same input always produces the same hash.
          $H(\text{``input''}) \to$ always the same output.

    \item \textbf{Avalanche Effect:} A single character change causes the entire
          hash to change completely and unpredictably. There is no gradual drift.

    \item \textbf{Fixed Size:} Regardless of whether the input is one word or
          one million pages, the output is always the same length (256 bits for SHA-256).

    \item \textbf{One-Way (Pre-image Resistance):} Given a hash $h$, it is
          computationally infeasible to find an input $x$ such that $H(x) = h$.
          You cannot reverse-engineer the original data from the fingerprint.
\end{enumerate}

\textbf{Example --- SHA-256 Avalanche Effect:}
\begin{center}
\begin{tabular}{l >{\ttfamily\small}l}
\toprule
\textbf{Input} & \textbf{SHA-256 Hash (truncated)} \\
\midrule
``Ali owes Shaheer 1000'' & a3f2c7d1...9e8b \\
``Ali owes Shaheer 1001'' & 7b84fa3c...2d07 \\
\bottomrule
\end{tabular}
\end{center}
One digit changed. The hash is completely unrecognisable.

\textbf{Why the one-way property is security-critical:}
\begin{itemize}
    \item \textbf{Privacy:} No one can read transaction details by inspecting a hash.
    \item \textbf{Forgery prevention:} An attacker who wants to secretly alter Block 2
          while keeping its hash identical (so Block 3 remains valid) would need to
          \emph{reverse} the hash function to engineer a collision. The one-way property
          makes this computationally infeasible, rendering such forgery impossible.
\end{itemize}

Formally, a secure hash function satisfies \textbf{collision resistance}: it is
infeasible to find two distinct inputs $x \neq x'$ such that $H(x) = H(x')$.

% ============================================================
\section{Phase 2: Consensus Mechanisms}
% ============================================================

\subsection{The Core Problem: Who Gets to Add Blocks?}

With thousands of strangers on the network, none of whom trust each other, a critical
question arises: who has the authority to append the next block? Appointing a gatekeeper
would recreate the centralisation problem. Allowing anyone to add freely would enable
spam and fraud.

The solution requires a mechanism with one key property:

\begin{center}
\textbf{Hard to produce. Trivially easy to verify.}
\end{center}

\begin{itemize}
    \item \textbf{Hard to produce} $\Rightarrow$ Adding a block costs real effort.
          A cheater cannot flood the network with fake blocks cheaply.
    \item \textbf{Easy to verify} $\Rightarrow$ Every node can confirm validity in
          milliseconds without trusting the producer.
\end{itemize}

\textbf{Sybil Attack:} If block creation were cheap, a single attacker could
masquerade as thousands of nodes and dominate the network. The cost of
production is the defence.

% ------------------------------------------------------------
\subsection{Proof of Work (PoW)}

Bitcoin's consensus mechanism. The ``puzzle'' is defined using the hash function:

\begin{quote}
    Find a number $n$ (called a \textbf{nonce}) such that:
    \[
        H(\text{block data} \;\|\; n) < T
    \]
    where $T$ is a target threshold, equivalent to requiring the hash output
    to begin with a certain number of leading zeros.
\end{quote}

\textbf{Example:}
\begin{verbatim}
    H(block_data + nonce) = 00000000003a7f2c...   ✓ valid
    H(block_data + nonce) = 7b84fa3c9d2e1a05...   ✗ invalid
\end{verbatim}

\textbf{Why this works:}
\begin{itemize}
    \item Because hashes are \textbf{one-way}, there is no shortcut. You cannot
          calculate the correct nonce --- you must guess, billions of times, until
          a valid hash appears by chance.
    \item Because hashes are \textbf{deterministic}, anyone can verify the answer
          by running the hash function exactly once.
\end{itemize}

The process of competing to find this nonce is called \textbf{mining}.

% ------------------------------------------------------------
\subsection{The Incentive Structure: Why Miners Are Honest}

A miner who wins (finds the valid nonce) earns:
\begin{enumerate}
    \item \textbf{Block reward} --- newly minted Bitcoin, created from nothing.
          This is the only mechanism by which new Bitcoin enters circulation.
    \item \textbf{Transaction fees} --- a small fee attached to every transaction
          by its sender.
\end{enumerate}

This creates a rational incentive table:

\begin{center}
\begin{tabular}{l l}
\toprule
\textbf{Miner Action} & \textbf{Outcome} \\
\midrule
Valid block + transactions included & Block reward + all fees \checkmark \\
Valid block + empty block            & Block reward only (fees forfeited) \\
Valid block + fraudulent transactions & Block rejected, zero reward, electricity wasted \texttimes \\
\bottomrule
\end{tabular}
\end{center}

\textbf{Key insight (Mechanism Design):} Satoshi did not ask miners to be honest.
He constructed the reward structure so that honesty is the \emph{most profitable
rational strategy}. The system enforces good behaviour through economics, not trust.

% ------------------------------------------------------------
\subsection{Difficulty Adjustment}

The target threshold $T$ is adjusted automatically every 2016 blocks
($\approx$ 2 weeks) so that the average time between blocks stays at
\textbf{10 minutes}, regardless of how much total computing power joins or
leaves the network.

More miners $\Rightarrow$ puzzles get harder. Fewer miners $\Rightarrow$ puzzles
get easier.

% ------------------------------------------------------------
\subsection{The 51\% Attack}

If a single entity controls more than half the network's total computing power
(hashrate), they can:
\begin{itemize}
    \item Consistently solve the puzzle faster than honest miners.
    \item Propose a secret alternative chain and eventually make it longer.
    \item Rewrite recent history (double-spend attacks).
\end{itemize}

This is called a \textbf{51\% attack}. Bitcoin's security therefore scales
directly with the total honest hashrate on the network. The larger and more
distributed the network, the more astronomically expensive such an attack becomes.

% ============================================================
\section{Phase 3: Proof of Stake \& Alternative Consensus}
% ============================================================

\subsection{The Problem with Proof of Work}

PoW is deliberately wasteful. Billions of hash guesses are computed solely to
produce a lottery ticket --- the actual computation creates no useful output.
Bitcoin's annual energy consumption rivals that of entire countries.

\textbf{Research question:} Can we replace ``burning electricity'' with something
else as proof of legitimate participation?

% ------------------------------------------------------------
\subsection{Proof of Stake (PoS)}

Instead of expending energy, validators \textbf{lock up} (stake) their coins as
collateral. The protocol pseudo-randomly selects a validator to propose the next
block, weighted by stake size:

\[
P(\text{selected}) \propto \text{stake}_i
\]

\textbf{Punishment --- Slashing:} If a validator is caught proposing fraudulent
blocks or signing contradictory blocks, the protocol automatically \textbf{destroys}
a portion of their staked coins. The attacker does not merely earn nothing --- they
lose what they committed.

\begin{center}
\begin{tabular}{l l l}
\toprule
\textbf{Property} & \textbf{Proof of Work} & \textbf{Proof of Stake} \\
\midrule
Right to add block     & Burn electricity          & Lock up coins \\
Punishment for cheating & Wasted electricity        & Slashing (coins destroyed) \\
Energy use             & Enormous                  & $\sim$99.9\% less \\
Hardware required      & Specialised ASICs         & Ordinary computer \\
Attack cost            & Buy 51\% of hashrate      & Buy 51\% of all staked coins \\
\bottomrule
\end{tabular}
\end{center}

\textbf{Wealth concentration problem:} A validator's selection probability scales
with stake. This means wealthier participants earn more rewards, compounding their
advantage --- a genuine ongoing tension in PoS system design.

Ethereum migrated from PoW to PoS in September 2022 (``The Merge'').

% ------------------------------------------------------------
\subsection{Delegated Proof of Stake (dPoS)}

\textbf{Motivation:} In plain PoS, hundreds of validators must coordinate to reach
consensus --- this is slow. dPoS introduces representative democracy to the blockchain.

\textbf{Mechanism:}
\begin{enumerate}
    \item Token holders \textbf{vote} to elect a small fixed set of trusted
          \textbf{delegates} (e.g.\ 21 in EOS).
    \item Only elected delegates produce blocks, taking turns in a
          round-robin schedule.
    \item Token holders can vote out a misbehaving delegate at any time.
\end{enumerate}

\textbf{Parliament analogy:} Token holders are citizens (voters); delegates are
elected MPs. Only MPs pass laws (add blocks), but citizens can replace them
at any election.

\begin{center}
\begin{tabular}{l l}
\toprule
\textbf{Property} & \textbf{dPoS Advantage} \\
\midrule
Speed           & Small delegate set $\Rightarrow$ fast consensus \\
Accountability  & Delegates voted out instantly for misbehaviour \\
Participation   & Anyone can vote regardless of stake size \\
\bottomrule
\end{tabular}
\end{center}

\textbf{Tradeoff:} Trust is concentrated in a small delegate group --- more
centralised than pure PoS or PoW.

\textbf{Two mechanisms that constrain a malicious delegate:}
\begin{enumerate}
    \item \textbf{Social (Voting):} Token holders immediately vote the delegate out,
          terminating all future block rewards and their position permanently.
    \item \textbf{Economic (Slashing):} Provable misbehaviour destroys the stake
          the delegate deposited to obtain their position. They paid to cheat
          and lose both the stake and the role.
\end{enumerate}

\textbf{CeDAR relevance:} The lab's paper on blockchain-enabled data sharing in
Connected Autonomous Vehicle (CAV) networks used dPoS with reputation schemes
(Multi-Weight Subjective Logic, beta, sigmoid) layered on top to detect malicious
nodes --- addressing the centralisation tradeoff through reputation scoring before
a node is even eligible for delegation.

% ============================================================
\section{Phase 4: Smart Contracts \& Ethereum}
% ============================================================

\subsection{The Core Idea}

Bitcoin's blockchain records \emph{who paid whom}. Ethereum extended this: what if the
blockchain could also \emph{run code}?

A \textbf{smart contract} is a program deployed permanently on the blockchain. No one
owns it, no one can stop it, and no one can modify it after deployment. When its
conditions are met, it executes automatically and irreversibly.

\textbf{Two fundamental limitations} (discovered through reasoning):
\begin{enumerate}
    \item \textbf{Incomplete contracts:} Code only handles what the programmer
          anticipated. Real-world edge cases (a draw, a disputed result, a force
          majeure) have no \texttt{else} clause. A judge can use judgment to fill
          legal gaps; a smart contract cannot. This is the \textbf{``Code is Law''
          problem.} Famous example: \emph{The DAO Hack (2016)} --- an attacker
          drained \$60M by exploiting the exact code as written. No rules were
          broken, because the code \emph{was} the rules.

    \item \textbf{The Oracle Problem:} A blockchain is a sealed, self-contained
          system. It cannot read the internet, watch a match, or sense the real
          world. Any contract that depends on external data requires a trusted
          data feed called an \textbf{oracle}. But trusting a single oracle
          re-introduces a centralised middleman. Solutions like \textbf{Chainlink}
          use decentralised oracle networks with their own staking/slashing
          mechanisms to make lying economically irrational.
\end{enumerate}

% ------------------------------------------------------------
\subsection{Storage and Execution: The EVM}

\textbf{Deployment:} A smart contract is deployed by sending a special transaction
whose payload contains the compiled contract bytecode. The network stores this code
permanently at a unique \textbf{contract address}. It cannot be deleted.

\textbf{Execution --- The Ethereum Virtual Machine (EVM):} Every node runs an
identical virtual machine. When a contract is called, every node executes the same
bytecode independently and must arrive at the same result. Disagreement results in
rejection. This mandates that smart contracts be \textbf{deterministic} --- given
the same inputs, they must always produce the same output. Non-deterministic
operations (random numbers, system time, external calls) are forbidden without
special handling.

% ------------------------------------------------------------
\subsection{Gas: Metered Computation}

\textbf{The halting problem:} A contract with an infinite loop would cause every
node to compute forever, halting the network.

\textbf{Solution --- Gas:} Every EVM instruction costs a fixed number of gas units.
The caller pays upfront. Execution stops when gas is exhausted.

\[
    \text{Transaction Cost (ETH)} = \text{Gas Used} \times \text{Gas Price}
\]

Gas serves three simultaneous functions:
\begin{enumerate}
    \item \textbf{Prevents infinite loops} --- execution halts when funds run out.
    \item \textbf{Prioritises transactions} --- validators include high-gas-price
          transactions first (rational economic choice).
    \item \textbf{Compensates validators} --- gas fees are the transaction fees
          validators collect per block.
\end{enumerate}

% ------------------------------------------------------------
\subsection{Access Control: The Foundational Security Pattern}

The single most important security pattern in smart contract development:

\begin{lstlisting}[language=Solidity, basicstyle=\ttfamily\small,
                   keywordstyle=\color{blue}, commentstyle=\color{gray}]
require(msg.sender == authorisedAddress, "Unauthorised");
\end{lstlisting}

\textbf{What it does:} Before executing any logic, verify \emph{who} is calling.
\texttt{msg.sender} is the cryptographically verified address of the caller ---
it cannot be forged. If the check fails, the entire transaction reverts as if it
never occurred.

\textbf{What it prevents:} Any attacker who tries to impersonate an authorised
party (e.g.\ an oracle, an owner, a regulator) to trigger fraudulent contract
execution --- such as calling \texttt{settle(true)} on a bet contract to falsely
claim a payout.

\textbf{Example --- Cricket Bet Contract:}

\begin{lstlisting}[language=Solidity, basicstyle=\ttfamily\small,
                   keywordstyle=\color{blue}, commentstyle=\color{gray}]
contract CricketBet {
    address payable public shaheer;
    address payable public ali;
    address public oracle;

    constructor(address payable _ali, address _oracle) payable {
        shaheer = payable(msg.sender); // deployer
        ali     = _ali;
        oracle  = _oracle;
    }

    function settle(bool pakistanWon) external {
        // Only the trusted oracle may call this
        require(msg.sender == oracle, "Only oracle can settle");
        if (pakistanWon) {
            shaheer.transfer(address(this).balance);
        } else {
            ali.transfer(address(this).balance);
        }
        // Note: draw case unhandled -- incomplete contract problem
    }
}
\end{lstlisting}

Key Solidity primitives:
\begin{itemize}
    \item \texttt{msg.sender} --- cryptographically verified caller address.
    \item \texttt{require()} --- guard clause; reverts entire transaction on failure.
    \item \texttt{transfer()} --- sends ETH atomically; no bank required.
    \item \texttt{address(this).balance} --- the ETH held by the contract itself.
\end{itemize}

% ------------------------------------------------------------
\subsection{CeDAR Relevance: Smart Contracts as Policy Enforcers}

CeDAR's core research encodes complex organisational access-control policies
(written in XACML) as smart contracts on Ethereum. The \texttt{require()} pattern
scales from a single address check to full \textbf{attribute-based access control
(ABAC)}: before granting a doctor access to a patient's record, the contract
verifies role, department, time-of-day, patient consent, and audit trail ---
all on-chain, all without a trusted administrator.

The gas cost of evaluating these policies on-chain is a primary research
challenge. CeDAR's work on \emph{policy composition} (merging multiple
organisations' policies into one composite smart contract) is specifically
motivated by minimising on-chain gas expenditure.

% ============================================================
\section{Phase 5: Cryptographic Identity --- Digital Signatures}
% ============================================================

\subsection{The Core Problem: Proving Authorship Without a Middleman}

When you send a transaction on Bitcoin, you broadcast a message to thousands of
strangers. How does every node on earth know that \emph{you} --- and not an
impersonator --- authorised it? There is no username, no password, no bank to
call. The answer is \textbf{public-key cryptography}.

% ------------------------------------------------------------
\subsection{The Magic Padlock: Asymmetric Keys}

A normal padlock uses one key to both lock and unlock (symmetric). The problem:
to give someone a locked box, you must first share the key --- but how do you
share it safely?

The solution is a padlock with a fundamentally different property:

\begin{quote}
    It comes with \textbf{two different keys}: a Private Key and a Public Key.\\
    What the Private Key \emph{locks}, only the Public Key can \emph{unlock}.\\
    Knowing the Public Key reveals \emph{nothing} about the Private Key.
\end{quote}

You keep the \textbf{Private Key} ($sk$) secret forever. You give the
\textbf{Public Key} ($pk$) to everyone --- post it publicly.

\textbf{How signing works:}
\begin{enumerate}
    \item You lock your message with your \emph{private} key and send it.
    \item Anyone uses your \emph{public} key to open it.
    \item The fact that the public key opened it \emph{proves} the private key
          locked it --- and only you have the private key.
\end{enumerate}

\textbf{The unforgeability guarantee:} Even if an attacker knows your public key
(everyone does), they cannot work backwards to find your private key. The
relationship between the two keys is one-way.

\begin{center}
\begin{tabular}{l l l}
\toprule
\textbf{Key} & \textbf{Kept by} & \textbf{Operation} \\
\midrule
Private key $sk$ & Owner only (secret) & \textbf{Sign} (lock) a message \\
Public key $pk$  & Everyone (public)   & \textbf{Verify} (unlock) a message \\
\bottomrule
\end{tabular}
\end{center}

% ------------------------------------------------------------
\subsection{The Mathematical Foundation: Modular Arithmetic}

All of this cryptography is built on \textbf{modular (clock) arithmetic}.

A 12-hour clock wraps around at 12:
\[
    15 \equiv 3 \pmod{12}
\]

\textbf{The Discrete Logarithm Problem:} Consider a clock with 7 positions.
Start at 1 and repeatedly add 3:
\[
    1 \;\to\; 4 \;\to\; 0 \;\to\; 3 \;\to\; 6 \;\to\; 2 \;\to\; 5 \;\to\; 1 \;\to\; \cdots
\]
The sequence jumps unpredictably. The question ``I added 3 exactly $k$ times
and landed on 5 --- what is $k$?'' has no algebraic shortcut. You must try
every possibility. For real cryptography the clock has $\sim 2^{256}$ positions:
finding $k$ requires $2^{128}$ operations. Computationally impossible.

% ------------------------------------------------------------
\subsection{Elliptic Curve Cryptography (ECC)}

Bitcoin and Ethereum use the curve \textbf{secp256k1}:
\[
    y^2 \equiv x^3 + 7 \pmod{p}
\]
Points on this curve form a group under \textbf{point addition} (draw a line
through two points, reflect the third intersection). Key generation:
\[
    P = k \cdot G = \underbrace{G + G + \cdots + G}_{k \text{ times}}
\]
\begin{itemize}
    \item $G$ --- public generator point (known to everyone)
    \item $k$ --- \textbf{private key} (random 256-bit integer, secret)
    \item $P$ --- \textbf{public key} (point on curve, shared openly)
\end{itemize}
Forward ($k \to P$): fast, $O(\log k)$ steps via repeated doubling.\\
Reverse ($P \to k$): the \textbf{ECDLP} --- requires $O(2^{128})$ operations.

% ------------------------------------------------------------
\subsection{Digital Signatures: ECDSA --- The Two Attacks and Their Defeat}

\textbf{Two attacks a forger could attempt:}
\begin{enumerate}
    \item \textbf{Forge without private key:} create a valid signature for a
          fraudulent transaction from scratch.
    \item \textbf{Replay attack:} copy a valid signature from one transaction
          and attach it to a different, fraudulent one.
\end{enumerate}

\textbf{Why hash before signing?} The message can be any length; the signing
math requires a fixed-size input. Hashing converts any message to exactly 256
bits: one hash, one signature, regardless of message size.

\textbf{Signing} (private key $k$, message hash $m$):
\begin{enumerate}
    \item $m = H(\text{transaction data})$
    \item Pick fresh random nonce $r$; compute $R = r \cdot G$
    \item $\displaystyle s = r^{-1}(m + k \cdot R_x) \pmod{n}$
    \item Output signature $\sigma = (R_x,\; s)$
\end{enumerate}

\textbf{Verification} (public key $P = kG$, no private key needed):
\[
    R' = s^{-1} \cdot m \cdot G + s^{-1} \cdot R_x \cdot P
\]
Valid iff $R'_x = R_x$.

\textbf{Why the algebra works} (substituting $P = kG$,
$s = r^{-1}(m + kR_x)$):
\[
    R' = \frac{r}{m+kR_x}\cdot m \cdot G + \frac{r}{m+kR_x} \cdot R_x \cdot kG
       = \frac{r(m + kR_x)}{m+kR_x} \cdot G = r \cdot G = R \;\checkmark
\]
The private key $k$ cancels completely --- proven without ever being revealed.

\textbf{Attack 1 defeated:} Forging $s$ requires $k$. Without it, the
verification equation produces $R'_x \neq R_x$.

\textbf{Attack 2 (replay) defeated:} $m$ is baked inside $s$. Attaching
$\sigma$ to a different message changes $m$, which changes $R'_x$. The
signature is \emph{message-bound} --- physically inseparable from the exact
transaction it was created for.

% ------------------------------------------------------------
\subsection{The Nonce Reuse Catastrophe}

\textbf{Critical rule:} The nonce $r$ must be freshly random for every signature.

\textbf{What happens if two signatures share the same $r$:}

Given $s_1 = r^{-1}(m_1 + k R_x)$ and $s_2 = r^{-1}(m_2 + k R_x)$:
\[
    s_1 - s_2 = r^{-1}(m_1 - m_2)
    \;\implies\;
    r = \frac{m_1 - m_2}{s_1 - s_2}
\]
And once $r$ is known:
\[
    k = \frac{s_1 \cdot r - m_1}{R_x}
\]
\textbf{The entire private key is exposed.} Every coin in the wallet is lost.
All future transactions can be forged. The damage is total and permanent.

\textbf{Real-world disaster --- Sony PlayStation 3 (2010):} Sony signed all
PS3 firmware with the \emph{same} nonce $r$. Researchers applied the above
algebra, recovered Sony's master private key, and could sign arbitrary software
as legitimate Sony firmware --- bypassing all security on every PS3 ever sold.

\textbf{Modern fix --- RFC 6979:} Rather than trusting a random number
generator, $r$ is derived \emph{deterministically} from the private key and
message hash:
\[
    r = \text{HMAC-SHA256}(k,\; m)
\]
This guarantees $r$ is unique per (key, message) pair without any randomness
that could fail or repeat.

% ------------------------------------------------------------
\subsection{Complete Signature Reference Table}

\begin{center}
\begin{tabular}{l l l}
\toprule
\textbf{Component} & \textbf{Value} & \textbf{Role} \\
\midrule
Private key $k$      & 256-bit secret          & Creates the signature \\
Public key $P = kG$  & Point on curve          & Verifies the signature \\
Message hash $m$     & $H(\text{transaction})$ & Binds signature to exact message \\
Nonce $R = rG$       & Fresh random point      & Prevents replay; must be unique \\
Signature $(R_x, s)$ & Two integers            & Proof of knowledge of $k$, bound to $m$ \\
\bottomrule
\end{tabular}
\end{center}

% ------------------------------------------------------------
\subsection{The Complete Bitcoin Transaction Flow}

\begin{verbatim}
YOUR WALLET
  Private key k  (secret, never leaves device)
       │  k · G   (one-way: ECDLP)
       ▼
  Public key P   (shared openly)
       │  Hash(P)  (one-way: SHA-256 + RIPEMD-160)
       ▼
  Address         (e.g. 1A2b3C... — what others send funds to)

SENDING A TRANSACTION
  1. Compose:  "Send 500 to Ali"
  2. Hash:      m = H("Send 500 to Ali")
  3. Sign:      σ = (Rx, s)  using private key k
  4. Broadcast: (message, σ, P) → all nodes

EVERY NODE VERIFIES INDEPENDENTLY
  Verify(P, m, σ) → true ✓  (no private key needed)
\end{verbatim}

\textbf{Security depth:} Two one-way barriers stand between your address and
your private key: ECDLP ($k \to P$) then hashing ($P \to$ address). An
attacker who sees only your address cannot recover $P$; even with $P$ they
cannot recover $k$.

% ============================================================
\section{Phase 6: DeFi, Tokens \& Tokenomics}
% ============================================================

\subsection{Why Custom Tokens? The Limits of Generic Currency}

ETH and BTC are \emph{generic} currencies --- they carry no rules about how they
must be spent. Paying for solar recycling in ETH means the recipient can spend
that ETH on anything. The money has no memory, no purpose, no constraints.

A \textbf{custom token} solves this by encoding spending rules \emph{into the
currency itself}:

\begin{itemize}
    \item A token can only be \textbf{minted} when a verified real-world event
          occurs (e.g.\ an oracle confirms a solar panel generated energy).
    \item A token can only be \textbf{burned} (redeemed) for a specific purpose
          (e.g.\ paying a certified recycler).
    \item The smart contract \textbf{refuses} all other uses.
\end{itemize}

This is called \textbf{programmable money} --- value that carries its own
spending rules. Generic ETH cannot do this because it has no awareness of
application-level context.

\textbf{CeDAR's RC-coin (Recycling Coin):} A custom token invented by CeDAR for
their solar waste research. Every unit of energy a panel generates triggers
minting of RC-coins. The accumulated coins represent the panel owner's recycling
obligation. At end-of-life, coins are burned to fund certified recycling. The
token is not merely money --- it is \emph{evidence of recycling obligation}
encoded as a transferable asset. No bank could enforce this; a smart contract
does it automatically and transparently.

% ------------------------------------------------------------
\subsection{The Three Token Types}

\begin{enumerate}
    \item \textbf{Fungible Tokens (ERC-20):} Every unit is identical and
          interchangeable, like banknotes. Used for currencies, governance votes,
          utility tokens. Standard interface ensures all wallets and exchanges
          can handle them without custom code.
          \textit{Examples: USDC, DAI, RC-coin, UNI.}

    \item \textbf{Non-Fungible Tokens (ERC-721):} Each token has a unique ID and
          distinct metadata. Used for provable ownership of specific assets: art,
          land titles, patient record access rights.

    \item \textbf{Stablecoins (a special ERC-20):} Pegged to a stable real-world
          value (\$1 USD). Two mechanisms:
          \begin{itemize}
              \item \textbf{Fiat-backed (USDC):} A company holds \$1 in a bank
                    for every USDC in existence. Simple but re-introduces a
                    centralised custodian.
              \item \textbf{Algorithmic / Crypto-collateralised (DAI):} No bank.
                    A smart contract maintains the peg using over-collateralisation
                    and liquidation mechanics --- pure DeFi.
          \end{itemize}
\end{enumerate}

% ------------------------------------------------------------
\subsection{Automated Market Makers (AMMs): Replacing Brokers with Math}

Traditional finance needs a \textbf{market maker} --- a bank or broker who
quotes buy/sell prices and takes a spread. DeFi replaces this with a single
mathematical invariant.

A \textbf{liquidity pool} holds two tokens. An \textbf{AMM} (e.g.\ Uniswap)
enforces the \textbf{constant product formula}:

\[
    x \cdot y = k
\]

where $x$ = quantity of Token A, $y$ = quantity of Token B, and $k$ is a
constant that never changes. No human sets prices. Price emerges from the ratio.

\textbf{Worked example.} Pool: 100 ETH, 10{,}000 RC-coin.
$k = 100 \times 10{,}000 = 1{,}000{,}000$.

You buy 10 ETH (removing it):
\[
    90 \cdot y = 1{,}000{,}000 \implies y = 11{,}111
\]
You paid $11{,}111 - 10{,}000 = 1{,}111$ RC-coin for 10 ETH. The ETH price rose
automatically (fewer ETH in pool). Supply and demand, encoded in one equation,
no intermediary.

% ------------------------------------------------------------
\subsection{Price Impact and the Sandwich Attack}

\textbf{Price impact (slippage):} A large trade moves the pool ratio
significantly, giving the buyer a worse average price. From $x \cdot y = k$:
removing a large $\Delta x$ from pool requires adding a disproportionately
large $\Delta y$. The formula punishes large trades automatically.

\textbf{Mathematical form of price impact:}
\[
    \Delta y = y \cdot \frac{\Delta x}{x - \Delta x}
\]
As $\Delta x \to x$, $\Delta y \to \infty$ --- you cannot drain a pool entirely.

\textbf{Sandwich Attack:} The mempool (pending transaction queue) is public. A
bot observing a large pending trade can:
\begin{enumerate}
    \item \textbf{Front-run:} buy Token A before the victim's transaction,
          pushing its price up.
    \item Let the \textbf{victim's transaction} execute at the now-inflated price.
    \item \textbf{Back-run:} immediately sell Token A at the higher price.
\end{enumerate}
The bot extracts value from the victim at zero risk. This is a form of
\textbf{Miner Extractable Value (MEV)}.

% ------------------------------------------------------------
\subsection{CeDAR's RC-Coin: DeFi Mechanics Applied to Solar Waste}

CeDAR's RC-coin system directly applies these tokenomics primitives:

\begin{center}
\begin{tabular}{l l}
\toprule
\textbf{DeFi Primitive} & \textbf{RC-coin Application} \\
\midrule
Token minting        & Triggered by energy-generation oracle \\
Token burning        & Redeemed to pay certified recycler \\
Stablecoin mechanism & RC-coin uses DeFi reserve mechanisms to \\
                     & stabilise recycling fund against ETH volatility \\
Smart contract rules & Enforce that RC-coin can \emph{only} fund recycling \\
On-chain audit trail & All recycling contributions are publicly verifiable \\
\bottomrule
\end{tabular}
\end{center}

Key finding from CeDAR's paper: the RC-coin scheme achieves a recycling
contribution of less than \$2/year/panel --- economically viable even for
prosumers in developing regions --- while maintaining a fully traceable,
tamper-proof contribution record on-chain.

% ============================================================
\section{Phase 7: CeDAR Research --- Deep Dive}
% ============================================================

% ------------------------------------------------------------
\subsection{Paper 1: Blockchain-Enabled Data Sharing in CAV Networks (ACNS 2023)}
% ------------------------------------------------------------

\textbf{Authors:} Ali Hussain Khan, Naveed Ul Hassan, Chuadhry Mujeeb Ahmed,
Zartash Afzal Uzmi.\\
\textbf{Venue:} Applied Cryptography and Network Security (ACNS) Workshops, 2023.

\subsubsection{The Real-World Problem}

A self-driving car cannot see around corners. Its safety depends on receiving
real-time data from other vehicles: obstacle locations, road conditions, sudden
braking events. But the network includes potentially hacked, spoofed, or
malicious nodes broadcasting false data to cause accidents or evade detection.

\textbf{The core threat:} A malicious node broadcasting false sensor data. Worse:
an \textbf{orchestrated attack} where multiple colluding malicious nodes all
broadcast the same lie and confirm each other's false reports, defeating
simple majority-vote systems.

\subsubsection{Why Blockchain? The Design Choice}

The paper uses blockchain as a \textbf{shared, tamper-proof reputation ledger}.
Every node's historical behaviour is recorded on-chain. Before any node's data
is trusted, its on-chain reputation score is checked.

\textbf{Consensus mechanism chosen: dPoS.} A small elected set of delegate nodes
produce blocks in round-robin. This achieves fast consensus --- critical in a
vehicular network where decisions must happen in milliseconds.

\textbf{Problem with vanilla dPoS in adversarial settings:} The delegate election
assumes the voting pool is trustworthy. If malicious nodes vote in malicious
delegates, the system collapses. The paper's solution is to layer a
\textbf{reputation gate} on top: a node cannot be eligible for delegation unless
its reputation score exceeds a threshold --- screening out adversaries
\emph{before} the election.

\subsubsection{Three Reputation Schemes Compared}

The paper evaluates three mathematical models for computing a node's reputation.

\textbf{1. Beta Reputation System}

Model each node's history as a sequence of binary outcomes: honest ($\alpha$ successes)
or dishonest ($\beta$ failures). Reputation is the expected value of the
Beta distribution:
\[
    R_{\text{beta}} = \frac{\alpha}{\alpha + \beta}
\]
Simple and interpretable. Weakness: all past evidence is weighted equally ---
recent misbehaviour counts no more than ancient history.

\textbf{2. Sigmoid Reputation}

Maps a scalar behaviour score $x$ to a bounded probability:
\[
    R_{\text{sig}} = \frac{1}{1 + e^{-x}}
\]
Smooth and bounded in $[0,1]$. Weakness: symmetric and single-dimensional ---
cannot distinguish between different \emph{types} of dishonest behaviour or
assign different severities to different failure modes.

\textbf{3. Multi-Weight Subjective Logic (MWSL)} --- the paper's primary proposal

Subjective Logic extends classical probability with an explicit \textbf{uncertainty
mass}: rather than asserting ``70\% honest,'' a node reports a
\emph{belief triplet} $(b, d, u)$ where $b$ = belief (honest), $d$ = disbelief
(dishonest), $u$ = uncertainty (insufficient data), satisfying $b + d + u = 1$.

The ``Multi-Weight'' extension evaluates behaviour across \textbf{multiple
independent dimensions}, each weighted by importance:
\begin{itemize}
    \item $w_1$: Sensor data accuracy (cross-verified against physical ground truth)
    \item $w_2$: Transmission timeliness (measured against a clock)
    \item $w_3$: GPS location consistency (cross-checked with base-station triangulation)
    \item $w_4$: Bandwidth/energy claims (verified against observed behaviour)
\end{itemize}

The combined reputation score:
\[
    R_{\text{MWSL}} = \sum_{i} w_i \cdot b_i, \qquad \sum_i w_i = 1
\]

\textbf{Why MWSL defeats the orchestrated attack:}

A majority-vote system fails against colluding nodes because 10 malicious nodes
can mutually confirm each other's false reports, appearing as consensus. MWSL is
immune because it \textbf{grounds reputation in verifiable external truth, not
just internal network agreement.} The colluding nodes can coordinate their lies,
but they cannot coordinate against a GPS base station, a physical clock, or
road-infrastructure sensor data that exists outside the network. Their false
reports contradict objective external references across multiple independent
dimensions simultaneously --- each contradiction reduces $b_i$, collapsing
$R_{\text{MWSL}}$ regardless of how many allies corroborate the lie.

\subsubsection{The Latency Model}

The paper does not merely claim blockchain is ``fast enough'' --- it formally
decomposes block-generation latency into four components:
\[
    T_{\text{block}} = T_{\text{comm}} + T_{\text{comp}} + T_{\text{prop}} + T_{\text{queue}}
\]

\begin{center}
\begin{tabular}{l l}
\toprule
\textbf{Component} & \textbf{Source of Delay} \\
\midrule
$T_{\text{comm}}$  & Wireless transmission of block data \\
$T_{\text{comp}}$  & Delegate computation / validation \\
$T_{\text{prop}}$  & Propagation of validated block to all nodes \\
$T_{\text{queue}}$ & Waiting time when channel is congested \\
\bottomrule
\end{tabular}
\end{center}

\subsubsection{Key Results}

The system is simulated under two attack models and three network generations:

\begin{center}
\begin{tabular}{l l l}
\toprule
\textbf{Network} & \textbf{Simple Attack} & \textbf{Orchestrated Attack} \\
\midrule
4G & Poor --- latency too high & Fails entirely \\
5G & Marginal                   & Poor \\
6G & Detected quickly           & \textbf{Detected within seconds} \\
\bottomrule
\end{tabular}
\end{center}

\textbf{Conclusion:} 4G and 5G networks have fundamentally unsuitable latency
for blockchain-based CAV consensus. 6G's sub-millisecond air interface collapses
$T_{\text{comm}}$ and $T_{\text{prop}}$, making MWSL-based dPoS viable for
real-time vehicular safety decisions.

\subsubsection{Interview Preparation: Key Discussion Points}

\begin{enumerate}
    \item \textbf{Why dPoS over PoW/PoS for CAVs?} Speed. A moving vehicle
          network cannot wait for mining or full-validator consensus. dPoS
          narrows the consensus set to a small elected group, minimising
          $T_{\text{comp}}$ and $T_{\text{prop}}$.

    \item \textbf{Why is reputation needed on top of dPoS?} dPoS assumes a
          trustworthy electorate. In adversarial CAV networks, the electorate
          itself may be compromised. Reputation gates the election, preventing
          malicious nodes from ever reaching delegate status.

    \item \textbf{What does MWSL add that beta/sigmoid cannot?}
          Multi-dimensional weighting. Different evidence types carry different
          danger levels. A lie about sensor accuracy (direct safety hazard) must
          penalise far more heavily than a delayed transmission. Beta and sigmoid
          are scalar --- they cannot express this hierarchy.

    \item \textbf{What is the open problem?} The paper validates in simulation.
          Real-world deployment requires standardised cross-verification
          infrastructure (physical reference nodes, base-station APIs) and
          a governance mechanism for setting weights $w_i$ --- who decides
          how much to penalise GPS inaccuracy vs.\ transmission delay?
\end{enumerate}

% ------------------------------------------------------------
\subsection{Paper 2: Blockchain-Based Auditable Access Control (ICDCS 2020 + TDSC 2024)}
% ------------------------------------------------------------

\textbf{Authors:} Ahmed Akhtar, Basit Shafiq, Jaideep Vaidya, Ayesha Afzal,
Shafay Shamail, Omer Rana (+ Masoud Barati for TDSC 2024).\\
\textbf{Venues:} IEEE ICDCS 2020; IEEE Transactions on Dependable and Secure
Computing (TDSC) 2024.

\subsubsection{The Real-World Problem}

A hospital, a research university, and a pharmaceutical company want to share
patient data. Each organisation has its own access control policy:
\begin{itemize}
    \item Hospital: ``Only cardiologists may access cardiac records during working hours.''
    \item University: ``Only researchers with active IRB approval may query patient data.''
    \item Pharma: ``Only employees with data-sharing agreement on file may read records.''
\end{itemize}

All three policies must be satisfied simultaneously before any access is granted.
Today this is enforced by a \textbf{central administrator} --- a trusted server
or person. Problems: single point of failure, can be bribed or hacked, no
tamper-proof audit trail.

\textbf{Shafiq's question:} Can we replace the central administrator with a smart
contract, while keeping on-chain execution cost economically viable?

\subsubsection{XACML: The Policy Language}

Organisations already write access policies in \textbf{XACML} (eXtensible Access
Control Markup Language), an XML-based standard for \textbf{Attribute-Based
Access Control (ABAC)}. Each rule tests \emph{attributes} of the requester:
role, department, clearance level, time-of-day, patient consent status. The
\texttt{require()} guard from Phase 4 is the primitive; XACML is its
industrial-strength formalisation.

\subsubsection{The Naïve Approach and Its Gas Cost Problem}

\textbf{Naïve design:} Deploy each organisation's policy as a separate smart
contract. For every access request, call all three contracts in sequence.

\textbf{Why this fails at scale:}

\begin{center}
\begin{tabular}{l l}
\toprule
\textbf{Approach} & \textbf{Per-Request On-Chain Cost} \\
\midrule
3 separate contracts (30 + 25 + 20 rules) & 3 separate calls, 75 rule evaluations \\
1 composite contract ($\sim$40 rules)     & 1 call, 40 rule evaluations \\
\bottomrule
\end{tabular}
\end{center}

At thousands of requests per day, three separate contracts cost approximately
3$\times$ the gas of a single composite. The on-chain cost makes the naïve
approach economically unviable at production scale.

\subsubsection{The Core Contribution: Policy Composition}

Instead of calling three contracts at runtime, \textbf{pre-merge} all three
policies into one composite policy \emph{at deployment time} (off-chain).
The composite enforces all three originals simultaneously --- one contract call,
one gas expenditure, per request.

\textbf{The deployment-time challenge:} With three policies of 30, 25, and 20
rules respectively, naïve composition must inspect up to
$30 \times 25 \times 20 = 15{,}000$ rule combinations to find overlaps,
redundancies, and contradictions. This combinatorial explosion is the
hard problem. It is paid once, off-chain, at composition time --- not on every
request. The key engineering insight is to \textbf{front-load complexity to
eliminate recurring on-chain cost}.

\textbf{Formal objective:}
\[
    \min |\mathcal{C}| \quad \text{subject to} \quad \forall i,\; \mathcal{C} \models \mathcal{P}_i
\]
Find the smallest composite policy $\mathcal{C}$ (fewest rules, minimum gas)
that logically implies every local policy $\mathcal{P}_i$.

\subsubsection{Two Formulations (ICDCS 2020 vs.\ TDSC 2024)}

\textbf{Paper 1 (ICDCS 2020) --- Attribute-Based Policies:}

Policies are conjunctions of attribute constraints. Composition is cast as a
\textbf{constraint optimisation problem}: find the minimum set of constraints
whose conjunction satisfies all local constraint sets. Redundant constraints
(subsumed by others) are eliminated. The result is deployed as a single
Ethereum smart contract with an Ethereum testnet prototype.

\textbf{Paper 2 (TDSC 2024) --- Event-Driven Policies:}

Some policies depend not on static attributes but on \textbf{sequences of events}:
``A researcher may access data \emph{only after} IRB approval has been recorded
\emph{AND} patient consent has been logged.''

These workflow policies are modelled as \textbf{finite automata} (state machines).
Composing $n$ automata into one minimal composite automaton reduces to the
\textbf{Weighted Set Cover problem}:

\begin{quote}
    Given a universe of required state transitions across all local automata,
    find the minimum-weight set of states in the composite that \emph{covers}
    all required transitions.
\end{quote}

Weighted Set Cover is NP-hard in general, but the XACML automaton structure
admits polynomial-time approximations with provable bounds --- which the paper
derives and validates empirically.

\subsubsection{Game-Theoretic Auditing}

Enforcement is solved by the composite smart contract. But \textbf{auditing} ---
verifying that the composite was derived correctly and that access was never
granted incorrectly --- remains expensive if done exhaustively on-chain.

Shafiq's solution: model auditing as a \textbf{game between a challenger
(auditor) and a defender (executing node)}. The challenger samples
probabilistically, not exhaustively. The mechanism design result:

\begin{quote}
    If the penalty for a detected violation exceeds the benefit of cheating,
    rational defenders never cheat --- regardless of the sampling probability.
\end{quote}

This is the same economic logic as Satoshi's miner incentive design (Phase 2),
now applied to policy auditing. A small probabilistic audit sample provides
the same deterrence as full auditing, at a fraction of the gas cost.

\subsubsection{Connection Map to Prior Phases}

\begin{center}
\begin{tabular}{l l}
\toprule
\textbf{Concept Learned} & \textbf{How It Appears in This Paper} \\
\midrule
\texttt{require()} guard (Phase 4) & Each XACML rule becomes one on-chain check \\
Gas cost (Phase 4)                 & Entire paper motivated by minimising gas \\
Smart contract immutability        & Composite deployed once, unalterable by any org \\
Oracle problem (Phase 4)           & Attributes (role, time) supplied via trusted oracles \\
Mechanism design / miners (Phase 2) & Auditing game: make cheating less profitable than honesty \\
\bottomrule
\end{tabular}
\end{center}

\subsubsection{Interview Preparation: Key Discussion Points}

\begin{enumerate}
    \item \textbf{What is the central cost tradeoff?} Front-load the
          combinatorial complexity of policy composition off-chain at
          deployment time (paid once). Eliminate recurring per-request gas by
          reducing 3 contract calls to 1. The offline cost is high
          ($O(n_1 \cdot n_2 \cdots n_k)$ combinations); the runtime saving
          compounds across millions of requests.

    \item \textbf{Why is Weighted Set Cover the right formulation for
          event-driven policies?} Workflow policies define required state
          transitions. A composite automaton must ``cover'' every transition
          required by every local automaton with the minimum number of states
          (each state = on-chain storage and computation = gas).

    \item \textbf{Why is full on-chain auditing impractical?} Every
          re-execution of a policy check costs gas. With millions of historical
          decisions, full re-audit is prohibitively expensive. Probabilistic
          sampling under a game-theoretic deterrence model achieves the same
          security guarantee at a fraction of the cost.

    \item \textbf{What is the open problem?} Policy \emph{updates}. When one
          organisation changes its policy, the composite must be recomposed
          and redeployed. The ablation studies in TDSC 2024 confirm that local
          policy changes correctly propagate through the composite, but
          redeployment costs gas and requires coordinating all participating
          organisations --- a governance problem the paper leaves open.
\end{enumerate}

\end{document}
